\subsection{Definitions}
Here are some definitions we will frequently use in the proof part. 
\begin{definition}
    [Even]
    An integer n is even if $n = 2a$ for some integer $a \in \mathbb{Z}$.
\end{definition}
\begin{definition}
    [Odd]
    An integer n is odd if $n = 2a + 1$ for some integer $a \in \mathbb{Z}$.
\end{definition}
\begin{definition}
    Two integers have the same parity if they are both even or they are both odd. Otherwise they have opposite parity. 
\end{definition}
\begin{definition}
    Suppose a and b are integers. We say that a divides b, written $a | b$, if $b = ac$ for some $c \in \mathbb{Z}$.
    In this case we also say that $a$ is a divisor of $b$, and that $b$ is a multiple of $a$. 
\end{definition}
\begin{definition}
    A number $n \in \mathbb{N}$ is prime if it has exactly two positive divisor, 1 and $n$. If $n$ has more than two positive divisors, it is called composite. 
\end{definition}
\begin{definition}
    [Greatest common divisor and Least common multiple]
    The greatest common divisor of integers a and b, denoted by $\gcd(a, b)$, is the largest integer that divides both a and b. 
    The least common multiple of non-zero integers a and b, denoted $lcm(a, b)$, is the smallest integer in $\mathbb{N}$ that is a multiple of both $a$ and $b$. 
\end{definition}

\newpage